% !PRINTED 44
\weberpage{44}
%<<CONT

\begin{equation}
\begin{aligned}
0&=B_0^{(0)}B_0^{(n)}-B_0^{(1)}B_1^{(n)}+B_0^{(2)}B_2^{(n)}-\cdots\pm B_0^{(n)}B_n^{(n)},\\
0&=-B_1^{(1)}B_1^{(n)}+B_1^{(2)}B_2^{(n)}-\cdots\pm B_1^{(n)}B_n^{(n)},\\
&\cdots\\
\pm 1&=\pm B_n^{(n)}B_n^{(n)},
\end{aligned}
\tag{6}
\end{equation}
and this system of formulas can also be written in the reversed order, using the relation $B_\nu^{(n)}=B_{n-\nu}^{(n)}$, as
\begin{equation}
\begin{aligned}
1&=B_0^{(n)}B_0^{(n)},\\
0&=B_1^{(n)}B_0^{(n)}-B_0^{(n-1)}B_1^{(n)},\\
0&=B_2^{(n)}B_0^{(n)}-B_1^{(n-1)}B_1^{(n)}+B_0^{(n-2)}B_2^{(n)},\\
&\cdots\\
0&=B_n^{(n)}B_0^{(n)}-B_{n-1}^{(n-1)}B_1^{(n)}+B_{n-2}^{(n-2)}B_2^{(n)}-\cdots\pm B_0^{(0)}B_n^{(n)},
\end{aligned}
\tag{7}
\end{equation}
or, in a summarizing notation,
\begin{equation}
 \sum_{\nu=0}^{\mu}(-1)^\nu B_{\mu-\nu}^{(n-\nu)}B_\nu^{(n)}
 =\begin{cases}
 1,&\mu=0,\\
 0,&\mu=1,2,\ldots,n.
 \end{cases}
\tag{8}
\end{equation}


\begin{center}
{\large \S.\ 9.}

\vspace{0.5em}
{\large\gesperrt{Interpolation.}}
\end{center}

We shall make an application of the last obtained formulae to a problem from the theory of interpolation.

\gesperrt{A whole rational function of degree $n$ is to be determined which, for $n+1$ given values of the variable $x$, has prescribed values.}

Thus the problem is the determination of the $n+1$ coefficients $a_0,a_1,\ldots,a_n$ in
\begin{equation}
 f(x)=a_0x^n+a_1x^{n-1}+\cdots+a_{n-1}x+a_n
\tag{1}
\end{equation}
from the $n+1$ equations
\begin{equation}
 f(\alpha_0)=A_0,\quad f(\alpha_1)=A_1,\quad \ldots,\quad f(\alpha_n)=A_n,
\tag{2}
\end{equation}
when $\alpha_0,\alpha_1,\ldots,\alpha_n$ are the given values of $x$, and $A_0,A_1,\ldots,A_n$ are the corresponding values of $f(x)$.

Equations (2) form a system of equations of the first degree for the unknowns $a_0,a_1,\ldots,a_n$, which in general can be solved by determinants, as we shall see in the next section. We shall later return to this problem,
%CONT>>
